\documentclass[10pt]{article}
\providecommand{\FIGDIR}{fig}
\newcommand{\DIGESTDATE}{17 August 2026}
% ===================== PM Research Watch - shared preamble =====================
\usepackage[T1]{fontenc}
\usepackage[utf8]{inputenc}
\usepackage{amsmath,amssymb}
\usepackage{tgpagella}
\usepackage{tgheros}
\renewcommand{\sfdefault}{qhv}
\renewcommand{\ttdefault}{lmtt}
\usepackage{microtype}
\usepackage[english]{babel}

\usepackage[a4paper,top=17mm,bottom=17mm,left=17mm,right=17mm,
            headsep=6mm,footskip=10mm]{geometry}
\usepackage{graphicx}
\usepackage{xcolor}
\usepackage{colortbl}
\usepackage{tikz}
\usetikzlibrary{calc}
\usepackage[most]{tcolorbox}
\usepackage{titlesec}
\usepackage{fancyhdr}
\usepackage{booktabs}
\usepackage{array}
\usepackage{longtable}
\usepackage{enumitem}
\usepackage{multicol}
\usepackage{ragged2e}
\usepackage[hidelinks]{hyperref}

% ------------------------------------------------------------------ palette
\definecolor{Ink}    {HTML}{14262E}
\definecolor{Deep}   {HTML}{0F4C5C}
\definecolor{Teal}   {HTML}{1F7A8C}
\definecolor{Sky}    {HTML}{4FA3C4}
\definecolor{Sage}   {HTML}{6FA287}
\definecolor{Amber}  {HTML}{E3A008}
\definecolor{Clay}   {HTML}{B5643C}
\definecolor{Coral}  {HTML}{D1495B}
\definecolor{Violet} {HTML}{6A5B9E}
\definecolor{Slate}  {HTML}{4A6572}
\definecolor{Paper}  {HTML}{FBF9F5}
\definecolor{Mist}   {HTML}{EEEAE2}
\definecolor{Rule}   {HTML}{D9D3C9}

\pagecolor{Paper}
\color{Ink}

\hypersetup{colorlinks=true, linkcolor=Deep, urlcolor=Teal, citecolor=Teal,
            pdftitle={PM Research Watch}, pdfauthor={Automated daily digest}}

% ------------------------------------------------------------------ headings
\titleformat{\section}
  {\sffamily\bfseries\large\color{Deep}}{}{0pt}
  {}[{\vspace{1.2mm}\color{Amber}\titlerule[1.1pt]}]
\titlespacing*{\section}{0pt}{5.5mm}{2.8mm}

\titleformat{\subsection}
  {\sffamily\bfseries\normalsize\color{Teal}}{}{0pt}{}
\titlespacing*{\subsection}{0pt}{3.6mm}{1.4mm}

\setlist[itemize]{leftmargin=4.2mm, itemsep=1.1mm, topsep=1.3mm, parsep=0pt,
                  label=\textcolor{Teal}{\footnotesize$\blacktriangleright$}}

% ------------------------------------------------------------------ header/footer
\pagestyle{fancy}
\fancyhf{}
\renewcommand{\headrulewidth}{0pt}
\renewcommand{\footrulewidth}{0pt}
\fancyhead[L]{\sffamily\scriptsize\color{Slate}PM RESEARCH WATCH\;\textcolor{Rule}{\textbar}\;\DIGESTDATE}
\fancyhead[R]{\sffamily\scriptsize\color{Slate}Particulate matter monitoring \& health effects}
\fancyfoot[C]{\sffamily\scriptsize\color{Slate}\thepage}

% ------------------------------------------------------------------ components
% metric tile
\newcommand{\metric}[3]{%
  \begin{tcolorbox}[enhanced, colback=white, colframe=Rule, boxrule=0.5pt,
    arc=1.6mm, left=2.2mm, right=2.2mm, top=1.8mm, bottom=1.8mm,
    borderline west={1.7mm}{0pt}{#3}, nobeforeafter, valign=center,
    height=17mm, width=\linewidth]
    \RaggedRight
    {\sffamily\bfseries\fontsize{17}{17}\selectfont\color{#3}#1}\\[0.5mm]
    {\sffamily\fontsize{6.5}{7.4}\selectfont\color{Slate}#2}
  \end{tcolorbox}}

% pull-quote / key insight box
\newtcolorbox{keybox}[1][]{enhanced, colback=Mist, colframe=Mist,
  arc=1.6mm, left=3.2mm, right=3.2mm, top=2.4mm, bottom=2.4mm,
  borderline west={1.7mm}{0pt}{Amber}, #1}

% subtopic band
\newcommand{\band}[2]{%
  \vspace{3.2mm}\noindent
  \begin{tcolorbox}[enhanced, colback=#2, colframe=#2, arc=1.2mm,
    left=2.6mm, right=2.6mm, top=1.1mm, bottom=1.1mm, boxrule=0pt,
    width=\linewidth, nobeforeafter]
    {\sffamily\bfseries\small\color{white}#1}
  \end{tcolorbox}\vspace{1.6mm}}

% one paper entry: {short cite}{title}{journal}{doi}{body}
\newcommand{\paperentry}[5]{%
  \par\vspace{1.9mm}\noindent
  {\sffamily\bfseries\footnotesize\color{Ink}#1}\;%
  {\sffamily\fontsize{7}{8}\selectfont\color{Slate}#3}\\[0.4mm]
  {\footnotesize\itshape\color{Deep}#2}\\[0.6mm]
  {\footnotesize #4}\\[0.3mm]
  {\fontsize{6.6}{7.6}\selectfont\color{Slate}\texttt{doi:}\,\href{https://doi.org/#5}{\texttt{\detokenize{#5}}}}
  \par}

% caption style
\newcommand{\figcap}[1]{\par\vspace{1.0mm}
  {\fontsize{7.2}{8.6}\selectfont\color{Slate}\textsf{\textbf{Fig.}}\; #1\par}\vspace{1.5mm}}

\providecommand{\permil}{\textperthousand}

\setlength{\parindent}{0pt}
\setlength{\parskip}{0pt}
\renewcommand{\arraystretch}{1.18}


\begin{document}

% ==================================================================== MASTHEAD
\thispagestyle{empty}
\begin{tikzpicture}[remember picture, overlay]
  \fill[Deep] (current page.north west) rectangle
      ($(current page.north east)+(0,-67mm)$);
  \fill[Amber] ($(current page.north west)+(0,-67mm)$) rectangle
      ($(current page.north east)+(0,-68.4mm)$);
  \foreach \x/\y/\r in {22/12/0.9, 41/26/1.5, 63/9/0.7, 88/31/1.1, 112/17/1.9,
                        138/29/0.8, 159/13/1.3, 178/54/1.0, 196/21/1.6,
                        30/44/0.6, 74/57/1.2, 125/48/0.9, 168/6/0.8, 191/59/0.7,
                        52/18/1.0, 100/60/0.6, 148/20/1.1, 66/49/1.4, 205/45/0.8}
    \fill[Sky, opacity=0.30] ($(current page.north west)+(\x mm,-\y mm)$) circle (\r mm);
\end{tikzpicture}

\vspace*{4mm}
{\sffamily\fontsize{7.6}{9}\selectfont\color{Sky}%
 AUTOMATED DAILY LITERATURE DIGEST\;\textbullet\;PhD RESEARCH WATCH\par}
\vspace{2.2mm}
{\sffamily\bfseries\fontsize{27}{29}\selectfont\color{white}%
 Particulate Matter\par}
{\sffamily\fontsize{27}{29}\selectfont\color{Sky}%
 Monitoring \& Health Effects\par}
\vspace{5.0mm}
{\sffamily\fontsize{8.4}{10}\selectfont\color{white}%
 Issue 2026-08-17\;\;\textcolor{Amber}{\textbar}\;\;Entry window 17 August 2026\;\;%
 \textcolor{Amber}{\textbar}\;\;12 records carried, 9 rejected\par}
\vspace{18mm}

% -------------------------------------------------------------- metric strip
\noindent
\begin{minipage}[t]{0.190\linewidth}\metric{12}{RECORDS CARRIED\\(9 REJECTED)}{Deep}\end{minipage}\hfill
\begin{minipage}[t]{0.190\linewidth}\metric{8}{SUBTOPIC CLUSTERS\\REPRESENTED}{Teal}\end{minipage}\hfill
\begin{minipage}[t]{0.190\linewidth}\metric{10}{EFFECT ESTIMATES\\WITH INTERVALS}{Sky}\end{minipage}\hfill
\begin{minipage}[t]{0.190\linewidth}\metric{6}{INSTRUMENTATION /\\EXPOSURE RECORDS}{Amber}\end{minipage}\hfill
\begin{minipage}[t]{0.190\linewidth}\metric{2}{METADATA-ONLY\\(LABELLED)}{Clay}\end{minipage}

\vspace{6mm}

% -------------------------------------------------------------- signal of day
\section*{\sffamily\bfseries\color{Deep}Signal of the day}
\vspace{-1.5mm}{\color{Amber}\rule{\linewidth}{1.1pt}}\vspace{2.5mm}

\begin{keybox}
{\sffamily\bfseries\small\color{Deep}The constituent ordering inverts against yesterday's
issue: here the secondary \emph{inorganic} ions lead and organic matter is the fraction
that collapses under adjustment.}\\[1.4mm]
\footnotesize
Zhao et al.\ (\textit{Chin Med J}) follow \textbf{20{,}117 pregnancies} in Nanjing with
residential PM$_{2.5}$ resolved into nitrate, sulfate, ammonium, black carbon and organic
matter, and estimate spontaneous-abortion odds per IQR by logistic regression with
propensity-score matching, WQS and BKMR. Total mass gives
\textbf{OR 1.18 (95\% CI 1.11--1.26)}; the ordering across constituents is
\textbf{NH$_4^+$ 1.19 (1.11--1.26)}, \textbf{NO$_3^-$ 1.18 (1.11--1.26)},
\textbf{SO$_4^{2-}$ 1.17 (1.10--1.25)}, then \textbf{OM 1.14 (1.07--1.21)} and
\textbf{BC 1.14 (1.08--1.22)}. After further adjustment for season of conception, OM
(1.08, 0.99--1.18) and BC (1.07, 0.97--1.18) go null while the three inorganic ions hold.
That is the exact reverse of the 16 August pairing, where Wu et al.\ attributed 54.1\% of a
prenatal joint effect to organic matter. The two are not in contradiction --- different
outcome, different cohort, different exposure product --- but taken together they say the
constituent that carries a prenatal association is not stable across studies, and the
seasonal-confounding structure is doing more work than either paper resolves.
\emph{For instrumentation the implication is sharper than yesterday's}: ammonium nitrate
and ammonium sulfate are the most hygroscopic fractions of the aerosol, so the constituents
that survive adjustment here are precisely the ones a nephelometric low-cost sensor
over-reads at high RH and precisely the ones a humidity-correction curve is fitted to
remove. A network calibrated to reproduce reference \emph{mass} is being tuned against the
signal this study says matters.
\end{keybox}

\vspace{3mm}

% -------------------------------------------------------------- exec brief
\section*{\sffamily\bfseries\color{Deep}What changed today}
\vspace{-1.5mm}{\color{Amber}\rule{\linewidth}{1.1pt}}\vspace{2.5mm}

\footnotesize
\begin{itemize}
\item \textbf{The greenness-modification claim in the same paper does not hold up to its
own numbers.} Zhao et al.\ report a significant PM$_{2.5}$$\times$NDVI-1000\,m interaction
and read it as greenness attenuating risk, but the quartile series is
\textbf{Q1 1.30 (1.14--1.48), Q2 1.02 (0.90--1.15), Q3 1.29 (1.14--1.46),
Q4 1.11 (0.97--1.26)} --- Q3 is statistically indistinguishable from Q1. That is
non-monotone, not a gradient, and the interaction $p$-value is being asked to carry a
dose--response story the strata do not support.

\item \textbf{A childhood-leukaemia signal arrives on a very thin exposure contrast.}
Choi et al.\ (\textit{Environ Res}) report \textbf{HR 1.40 (1.04--1.90) per 5\,$\mu$g/m$^3$}
across 384{,}606 Korean children, but the cohort mean is
\textbf{29.6\,$\mu$g/m$^3$ with an SD of only 3.1} and exposure is assigned at
district level from an ML ensemble. A 5\,$\mu$g/m$^3$ increment is more than 1.5 SD of the
available between-district spread, so the estimate is an extrapolation across most of its
own range on 272 events.

\item \textbf{Proteomic signatures are proposed as a replacement for the ambient
concentration --- and cannot be checked here.} Chen et al.\ (\textit{Ecotox Environ Saf})
screen $\sim$2{,}900 UK Biobank plasma proteins against seven pollutants, build elastic-net
signatures, and report they out-predict external exposure metrics for incident AF, with
ANGPT2, GDF15 and PLAUR as mediators and \textbf{1.15--1.25 years of AF-free life lost} in
the top-risk half. No hazard ratio or interval appears in the abstract, so the record
carries \emph{no} estimate into today's forest plot despite being the day's largest study.

\item \textbf{One directly actionable siting result.} Ianiri et al.\
(\textit{Environ Pollut}) find street-level deposition fluxes roughly \textbf{1.6$\times$
roof-level} (42 vs 26\,mg\,m$^{-2}$\,d$^{-1}$ annual SPM) but \emph{no} significant
street--roof difference once pollutant loads are normalised to collected SPM
(Wilcoxon, $p>0.05$). Sampling height changes how much dust you collect, not what the dust
is made of --- which is the assumption most rooftop-versus-breathing-zone network designs
implicitly make and rarely test.

\item \textbf{Three optical-property records in one day, and two ship without an abstract.}
The Halo Doppler lidar depolarisation work is complete and quantitative; the AS\&T humidified
CAPS-PM$_{\mathrm{SSA}}$ paper and the \textit{J Aerosol Sci} refractive-index inversion are
carried on title and author only, because Taylor \& Francis and Elsevier both withheld the
abstract and Semantic Scholar was rate-limited for the whole run. Both are labelled
\texttt{Metadata only} rather than pooled with characterised designs. The nine rejections,
including two scope calls on PM-adjacent reviews, are itemised in the provenance note.
\end{itemize}

\clearpage

% ==================================================================== ANALYTICS
\section*{\sffamily\bfseries\color{Deep}Corpus analytics}
\vspace{-1.5mm}{\color{Amber}\rule{\linewidth}{1.1pt}}\vspace{3mm}

\noindent\includegraphics[width=\linewidth]{\FIGDIR/f1_subtopics.png}
\figcap{Subtopic assignment of the 12 carried records, single-label by dominant
contribution. Eight of ten clusters are occupied and no cluster holds more than three, which
is 0.67 clusters per record --- \textbf{the most dispersed issue of the watch to date}, ahead
of 14 August (0.64) and 31 July (0.64), computed over all 23 issues.
\textit{Sensing, forecasting \& instrumentation} and
\textit{Exposure assessment \& modelling} tie at 3 each; six clusters hold exactly one
record. \textit{Neuro / mental health} is empty for the twelfth time in twenty-three issues
and \textit{Occupational \& indoor} for the sixth, both counts computed from
\texttt{state/metrics.csv} rather than recalled. Twelve records ties 9 August as the
fourth-smallest issue of the watch, behind 10 (5 and 16 August) and 11 (31 July).}

\vspace{2mm}
\noindent\includegraphics[width=\linewidth]{\FIGDIR/f2_design_endpoint.png}
\figcap{Left --- study architecture. Observational cohorts lead at 3 (Nanjing pregnancies,
Korean leukaemia, UK Biobank proteomics), matched by 3 measurement campaigns (school metals,
Rome deposition, Halo lidar). \textbf{\textit{Metadata only} takes 2 of 12} and is shown as
its own slice by design: pooling an uncharacterised record with a characterised one is how
the architecture panel was quietly wrong before this category existed. Right --- health
endpoint. \textbf{No health endpoint} is again the largest bar at 6 of 12 (50\%), sitting
mid-band in the 41--67\% range of the previous ten issues (7--16 August: 54, 64, 67, 56, 58,
46, 41, 43, 43, 60\%).}

\vspace{2mm}
\noindent\includegraphics[width=\linewidth]{\FIGDIR/f3_metric_geography.png}
\figcap{Left --- particle metric. PM$_{2.5}$-only accounts for 5 records and PM$_{2.5}$
with PM$_{10}$ jointly for 2; \textbf{optical properties (absorption/scattering) take 3},
the whole instrumentation block, and one record each covers ultrafine carbon black and
settleable dust. Right --- geography. Europe leads at 4 (Sweden, Italy, Finland, UK).
\textbf{The \textit{Global / multi-region} slice of 4 should be read as an absence, not as
reach}: only the GLENS preprint is genuinely global; the other three are a murine
laboratory study, a bench instrument paper and a metadata-only algorithm record with no
stated field site. South Asia's single record is the Pakistan satellite study, the first
since 15 August; it is also the only record today from a region carrying a top-decile
attributable burden.}

\vspace{2mm}
\noindent\includegraphics[width=\linewidth]{\FIGDIR/f6_lifecourse.png}
\figcap{Life-course windows addressed. Counts are non-exclusive and describe this issue
only. Preconception/in-utero, infancy, adolescence and working age each hold 2; the 65+
window holds 1 and only by inference, since UK Biobank participants age into it across
13.5 years of follow-up rather than being recruited there. The leukaemia cohort contributes
to three consecutive windows because it enrols at birth and follows past age 18, which is
the design feature that distinguishes it from the rest of the paediatric literature.}

\vspace{2mm}
\noindent\includegraphics[width=\linewidth]{\FIGDIR/f4_heatmap.png}
\figcap{Subtopic $\times$ study-architecture cross-tabulation. Twelve records occupy 9 of
the 80 available cells and no cell holds more than 2. The silo persists in the same shape as
the last three issues: every \textit{Sensing} and \textit{Exposure} record sits in a
measurement, modelling or metadata-only column and every human-endpoint record in an
observational one, with no cell bridging them. Note that the two \textit{Metadata only}
entries sit entirely inside \textit{Sensing} --- the publishers that withhold abstracts on
this watch are the aerosol-instrumentation publishers, so the gap is systematic rather than
incidental.}

\vspace{2mm}
\noindent\includegraphics[width=\linewidth]{\FIGDIR/f5_forest.png}
\figcap{The 10 interval estimates carried today, on a common logarithmic axis.
\textbf{Eight of the ten come from one study and are not independent.} The six Zhao
constituent estimates are IQR-scaled contrasts on correlated fractions of the same aerosol
in the same 20{,}117 pregnancies, so their \emph{ordering} is informative and their spread is
not a between-study spread; the two greenness strata are subsets of that same cohort and
share person-time with all six. Only Tornevi (per IQR 3.3\,$\mu$g/m$^3$, 1-day lag,
case-crossover) and Choi (per 5\,$\mu$g/m$^3$, annual mean, cohort) are independent of it and
of each other, and their exposure contrasts differ by roughly an order of magnitude in both
averaging time and increment. Read within family, never across.}

\clearpage

% ==================================================================== DIGEST
\section*{\sffamily\bfseries\color{Deep}Paper-level digest}
\vspace{-1.5mm}{\color{Amber}\rule{\linewidth}{1.1pt}}\vspace{1mm}

{\fontsize{7.4}{9}\selectfont\color{Slate}Ordered by cluster. Each entry gives the
methodological core and the finding that matters, not an abstract paraphrase.\par}

\band{Cardiovascular \& metabolic\hfill 1 record}{Teal}

\paperentry{Chen et al.\ 2026}
{Pollutant-specific plasma proteomic signatures and atrial-fibrillation-free life lost}
{Ecotoxicol Environ Saf}
{UK Biobank, $\sim$2{,}900 Olink plasma proteins screened against long-term PM$_{2.5}$,
PM$_{10}$, NO$_2$, NO$_x$, SO$_2$, benzene and O$_3$ under Bonferroni correction, then
elastic-net regularisation to build one signature per pollutant; incident AF over a median
\textbf{13.53\,y} by Cox regression. Signatures out-predicted the external exposure metrics
and stratified risk after adjusting for ambient concentration; high-risk participants lost
an estimated \textbf{1.15--1.25 AF-free years}. High-throughput mediation nominates
\textbf{ANGPT2, GDF15 and PLAUR}. The weakness is structural, not analytic: a signature
trained on a modelled exposure inherits that model's error and its spatial confounding, so
"out-performs the exposure metric" may mean "encodes residence better". No HR or interval
is reported, so nothing enters today's forest plot.}
{10.1016/j.ecoenv.2026.120666}

\band{Reproductive \& developmental\hfill 1 record}{Sky}

\paperentry{Zhao et al.\ 2026}
{PM$_{2.5}$ constituents and spontaneous abortion, modified by residential greenness}
{Chin Med J (Engl)}
{20{,}117 pregnancies in Nanjing with residential PM$_{2.5}$ speciated into NO$_3^-$,
SO$_4^{2-}$, NH$_4^+$, BC and OM; logistic regression with propensity-score matching plus
two-pollutant models, WQS and BKMR for the joint effect. Per IQR: total mass
\textbf{OR 1.18 (1.11--1.26)}, ammonium \textbf{1.19 (1.11--1.26)} highest,
BC and OM lowest at 1.14 --- and OM and BC alone lose significance once season of conception
is added. Greenness (NDVI-1000\,m) interacts significantly but non-monotonically
(Q1 1.30, Q2 1.02, Q3 1.29, Q4 1.11), which undercuts the attenuation reading. Depression
and anxiety scores predicted SAB independently with no interaction. Constituent
correlations are high, so the ordering is suggestive rather than separable.}
{10.1097/CM9.0000000000004259}

\band{Respiratory \& allergic\hfill 1 record}{Sage}

\paperentry{Tornevi et al.\ 2026}
{Short-term urban-background PM and the timing of confirmed COVID-19, three Swedish cities}
{BMJ Open Respir Res}
{Case-crossover on \textbf{3{,}999} PCR-confirmed cases from a prospective adult cohort in
Stockholm, Gothenburg and Uppsala, 7 Mar 2020 -- 26 Jan 2022, with daily PM$_{2.5}$,
PM$_{10}$ and NO$_2$ from urban-background stations and distributed-lag (and non-linear)
models adjusted for meteorology. \textbf{PM$_{2.5}$ at lag 0: OR 1.084
(1.011--1.162) per IQR 3.3\,$\mu$g/m$^3$}, consistent across all three cities; PM$_{10}$
most consistent at lag 3; NO$_2$ lag 0 only after dropping the final wave. The authors
themselves read the same-day lag as implausible for infection and propose symptom
exacerbation prompting testing --- i.e.\ the outcome is a \emph{test date}, not an infection
date, and the estimate is at least partly a health-seeking-behaviour effect.}
{10.1136/bmjresp-2025-003279}

\band{Mechanistic toxicology\hfill 1 record}{Violet}

\paperentry{Li et al.\ 2026}
{Size-resolved carbon-black nanoparticle effects on the maternal--fetal lung axis}
{FASEB J}
{Pregnant mice exposed to 30\,nm or 120\,nm carbon black under an explicitly high-dose
mechanistic protocol, with H\&E/Masson histology, dark-field hyperspectral imaging for
particle localisation and chemical-isotope-labelling LC-MS metabolomics. Both sizes produced
maternal inflammatory infiltration and fibrosis, \textbf{120\,nm worse than 30\,nm};
transplacental transfer was imaged directly in fetal lung. Fetal lungs showed no histological
injury but raised oxidative and inflammatory markers, with size-specific metabolic
signatures --- 30\,nm on histidine metabolism, 120\,nm on butanoate and arachidonic acid.
Two problems: the dose is not environmentally referenced, and the larger particle causing
more damage inverts the surface-area expectation without a number- versus mass-dose
accounting that would explain it.}
{10.1096/fj.202600065RR}

\band{Exposure assessment \& modelling\hfill 3 records}{Amber}

\paperentry{Moryani et al.\ 2026}
{PM$_{2.5}$-bound Cu, Cd, Ni, Pb and Zn across 20 schools in two contrasting South China cities}
{J Appl Toxicol}
{Paired summer and winter sampling at 20 schools in Guangzhou and Maoming, ICP-MS for five
metals, with geoaccumulation index, enrichment factor and PCA source apportionment, then
inhalation risk assessment for children and adults. Cu and Cd carried the highest enrichment
(EF 27.72 and 27.32 in Guangzhou), winter exceeded summer, and PCA pointed to traffic,
industrial combustion and urban emissions. \textbf{All hazard quotients and carcinogenic
risks fell below accepted thresholds}, with Pb (4.85$\times$10$^{-3}$) and Cu
(3.27$\times$10$^{-3}$) leading non-carcinogenic risk and Cd (7.30$\times$10$^{-6}$)
carcinogenic. Fieldwork dates from 2018, so it characterises an emissions regime seven years
old; the null risk result is the honest headline and is under-stated in the abstract.}
{10.1002/jat.70398}

\paperentry{Ianiri et al.\ 2026}
{Street- versus roof-level deposition fluxes of POPs and metals, normalised to sedimentable PM}
{Environ Pollut}
{Thirteen monthly deposimeter samples at two heights on one Rome building, Aug 2023 --
Aug 2024, with PCDD/Fs, PCBs and PAHs by GC-HRMS and metals by ICP-MS, reported both as
fluxes and normalised to collected sedimentable PM. \textbf{Annual SPM flux 42 vs
26\,mg\,m$^{-2}$\,d$^{-1}$} street versus roof, but a Wilcoxon signed-rank test found
\textbf{no significant street--roof difference in SPM-normalised concentrations
($p>0.05$)}: height changes deposited quantity, not deposited composition. Seasonality was
significant for PCDD/Fs, PCBs, Cr, As, Ni and Sb (winter peaks) and for PAHs
(summer/autumn). Single building, single city, $n=13$ per height --- a design that supports
the composition claim far better than the flux ratio.}
{10.1016/j.envpol.2026.128965}

\paperentry{Awan et al.\ 2026}
{BRIQ-PM$_{2.5}$: quantile-mapping attribution of industrial excess PM$_{2.5}$ across Pakistan}
{Aerosol Air Qual Res}
{Proposes estimating industrial impact as the residual between an observed satellite
PM$_{2.5}$ distribution and the same distribution quantile-mapped onto a province-specific
"low-industry" baseline city, then interpolating the residual by variogram modelling and
universal kriging across four provinces. Reported validation is distributional --- histogram
and ECDF alignment, improved RMSE after correction. \textbf{The attribution is definitional,
not estimated}: whatever separates a target city from its chosen baseline is labelled
industrial, so meteorology, dust, biomass burning and population density all load onto the
same term, and no ground-monitor comparison or sensitivity analysis over baseline choice is
reported. Useful as a hotspot screen; not an emissions attribution.}
{10.1007/s44408-026-00160-z}

\band{Sensing, forecasting \& instrumentation\hfill 3 records}{Amber}

\paperentry{Filioglou et al.\ 2026}
{Particle linear depolarisation ratio at 1565\,nm from Halo Doppler lidar}
{Atmos Chem Phys}
{Assesses whether a wind-profiling Halo Photonics StreamLine Doppler lidar can type aerosol
at 1565\,nm, a wavelength where depolarisation is rarely characterised, across three
tropospheric case studies. \textbf{Extremely fresh smoke gave 0.017 $\pm$ 0.004}, aged
long-range-transported smoke marginally higher, and \textbf{volcanic ash layers 0.45
$\pm$ 0.01} --- a separation wide enough to type layers unambiguously. Smoke mass
concentration retrieved from the 1565\,nm signal agreed with concurrent in-situ
observations. The value here is that an instrument already deployed for wind now also
delivers aerosol typing and a mass estimate on the same beam; three case studies is thin
for a conversion factor, and the aged-smoke result is qualitative.}
{10.5194/acp-26-11525-2026}

\paperentry{Dal Porto et al.\ 2026}
{Humidified CAPS-PM$_{\mathrm{SSA}}$ with a revised truncation-correction methodology}
{Aerosol Sci Technol}
{\textbf{Metadata only.} Taylor \& Francis withheld the abstract on this run and Semantic
Scholar was rate-limited throughout, so nothing below the title, authors and journal is
verified and no result is asserted. On the registered title alone the paper describes a
humidified variant of the cavity-attenuated phase-shift single-scattering-albedo monitor
together with an improved method for the scattering truncation correction --- the two
quantities that dominate uncertainty when an SSA measurement is used as a reference for
hygroscopic growth. Relevant enough to the humidity-correction problem to carry rather than
defer; flagged for abstract retrieval on the next run.}
{10.1080/02786826.2026.2711663}

\paperentry{Zhao et al.\ 2026}
{Monte-Carlo inversion of aerosol complex refractive index from multi-wavelength lidar}
{J Aerosol Sci}
{\textbf{Metadata only.} Elsevier returned no abstract through Crossref and ScienceDirect
did not serve the record; no finding is asserted. The registered title describes an improved
Monte-Carlo sampling algorithm that inverts the aerosol complex refractive index from
multi-wavelength lidar signals combined with an independently measured particle size
distribution. Carried at tier C because refractive index is the parameter that converts an
optical measurement into a mass estimate, and every low-cost nephelometric sensor assumes one
implicitly; whether this inversion is well-posed cannot be judged from the title.}
{10.1016/j.jaerosci.2026.106883}

\band{Burden, policy \& mitigation\hfill 1 record}{Coral}

\paperentry{Xia et al.\ 2026 \textnormal{(preprint)}}
{Surface PM$_{2.5}$ mass and composition under stratospheric sulfate geoengineering}
{arXiv}
{Uses the GLENS ensemble (CESM1, stratospheric SO$_2$ injection balancing RCP8.5 forcing) to
ask what solar geoengineering does to surface PM$_{2.5}$. Global mean PM$_{2.5}$
\emph{falls} relative to RCP8.5, but the reduction is carried entirely by \textbf{less dust}
(wetter soils, higher leaf-area index over deserts) and \textbf{less sea salt} (slower
winds). Excluding those two, PM$_{2.5}$ \emph{rises}, driven by aerosol-phase secondary
organic aerosol partitioning in a cooler atmosphere. Notably, the injected SO$_2$ does
\textbf{not} raise surface sulfate. Single model, single ensemble, an RCP8.5 baseline now far
from plausible, and SOA gas precursors are prescribed --- so the SOA increase is partly an
artefact of that setup, as the authors acknowledge. Preprint; not peer reviewed.}
{10.48550/arXiv.2608.14928}

\band{Other clinical endpoints\hfill 1 record}{Slate}

\paperentry{Choi et al.\ 2026}
{Time-varying long-term PM$_{2.5}$ and leukaemia from childhood into early adulthood}
{Environ Res}
{National Health Insurance Service--National Sample Cohort, 2002--2020:
\textbf{384{,}606} people aged 0--18 at enrolment, \textbf{272 incident leukaemias}
(83 diagnosed after age 18), annual PM$_{2.5}$ assigned to all 226 inland districts from an
ML ensemble, Cox models with time-varying exposure and confounders plus two-pollutant
sensitivity against NO$_2$ and O$_3$. \textbf{HR 1.40 (1.04--1.90) per 5\,$\mu$g/m$^3$},
stable when censored at 18. The exposure contrast is the problem: cohort mean
\textbf{29.6 $\pm$ 3.1\,$\mu$g/m$^3$}, so a 5\,$\mu$g/m$^3$ increment exceeds 1.5 SD of the
between-district spread and the estimate extrapolates beyond most of its own support, on
272 events with district-level assignment and no residential-mobility handling.}
{10.1016/j.envres.2026.125497}

\clearpage

% ==================================================================== TABLE
\section*{\sffamily\bfseries\color{Deep}Corpus register}
\vspace{-1.5mm}{\color{Amber}\rule{\linewidth}{1.1pt}}\vspace{2.5mm}

{\fontsize{6.7}{8.4}\selectfont
\setlength{\tabcolsep}{2.4pt}
\begin{longtable}{@{}>{\RaggedRight}p{25mm}>{\RaggedRight}p{41mm}>{\RaggedRight}p{27mm}%
>{\RaggedRight}p{23mm}>{\RaggedRight}p{18mm}>{\RaggedRight}p{15mm}@{}}
\toprule
\textsf{\textbf{First author}} & \textsf{\textbf{Focus}} & \textsf{\textbf{Journal}} &
\textsf{\textbf{Design}} & \textsf{\textbf{Metric}} & \textsf{\textbf{Region}}\\
\midrule
\endfirsthead
\toprule
\textsf{\textbf{First author}} & \textsf{\textbf{Focus}} & \textsf{\textbf{Journal}} &
\textsf{\textbf{Design}} & \textsf{\textbf{Metric}} & \textsf{\textbf{Region}}\\
\midrule
\endhead
\multicolumn{6}{@{}l@{}}{\cellcolor{Mist}\textsf{\textbf{\textcolor{Deep}{Neuro / mental health}}}}\\[0.6mm]
\href{https://doi.org/10.1001/jamanetworkopen.2026.25015}{Dales} & Wildfire-sourced PM$_{2.5}$ and ED visits for migraine and other primary headache & \textit{JAMA Netw Open} & Case-crossover & PM$_{2.5}$ (wildfire) & Canada\\
\href{https://doi.org/10.1093/ageing/afag215}{Glover} & Umbrella meta-review of urban environmental exposures and cognitive health & \textit{Age Ageing} & Umbrella review / meta-analysis & PM$_{2.5}$ & Global\\
\href{https://doi.org/10.1016/j.nsa.2026.107019}{Goossens} & In utero air pollution exposure and pre-/postnatal brain development & \textit{Neurosci Appl} & Narrative review & PM$_{2.5}$ / PM$_{10}$ / PM$_{\mathrm{coarse}}$ & Global\\
\href{https://doi.org/10.1016/j.isci.2026.116617}{Ma} & PM$_{2.5}$-induced neurotoxicity and depression: exposure risks and mechanisms & \textit{iScience} & Narrative review & PM$_{2.5}$ + components & Global\\
\href{https://doi.org/10.3760/cma.j.cn112150-20250804-00752}{Pan} & Maternal PM exposure and offspring neurodevelopment: effects and mechanisms & \textit{Zhonghua Yu Fang Yi Xue Za Zhi} & Narrative review & PM (bulk) & China\\
\addlinespace[0.9mm]
\multicolumn{6}{@{}l@{}}{\cellcolor{Mist}\textsf{\textbf{\textcolor{Teal}{Cardiovascular \& metabolic}}}}\\[0.6mm]
\href{https://doi.org/10.3760/cma.j.cn112338-20251120-00843}{Huang} & PM$_{2.5}$ and dyslipidaemia in rural Chinese children; dietary effect modification & \textit{Zhonghua Liu Xing Bing Xue Za Zhi} & Cross-sectional (multistage) & PM$_{2.5}$ & China\\
\href{https://doi.org/10.32394/pe/218966}{Kazmierska} & Air pollution and cardiovascular disease: review of current data & \textit{Przegl Epidemiol} & Narrative review & PM$_{2.5}$ / PM$_{10}$ & Global\\
\href{https://doi.org/10.1007/s00125-026-06814-2}{Xi} & Air pollution, built environment and temperature vs. prediabetes phenotypes & \textit{Diabetologia} & Prospective cohort & PM$_{2.5}$ & Germany\\
\href{https://doi.org/10.1161/JAHA.125.046917}{Yang} & Short-term PM$_{2.5}$ and hospitalisation/ED use in atrial fibrillation & \textit{J Am Heart Assoc} & Case-crossover & PM$_{2.5}$ (wildfire-influenced) & USA\\
\addlinespace[0.9mm]
\multicolumn{6}{@{}l@{}}{\cellcolor{Mist}\textsf{\textbf{\textcolor{Sky}{Reproductive \& developmental}}}}\\[0.6mm]
\href{https://doi.org/10.1016/j.jes.2025.10.043}{Chen} & Prenatal air pollution, fetal growth trajectory phenotypes and child neurodevelopment & \textit{J Environ Sci (China)} & Prospective cohort & PM$_{2.5}$ / PM$_{10}$ & China\\
\href{https://doi.org/10.1016/j.ecoenv.2026.120566}{Huang} & Joint prenatal bisphenol + air-pollutant mixture exposure and preschool IQ & \textit{Ecotoxicol Environ Saf} & Prospective cohort (mixtures) & PM$_{2.5}$ / PM$_{10}$ & China\\
\href{https://doi.org/10.1007/s40572-026-00552-8}{Peterson} & Ambient air pollution, PAHs and pubertal development: critical review & \textit{Curr Environ Health Rep} & Critical review & PM$_{2.5}$ components & Global\\
\href{https://doi.org/10.1016/j.envres.2026.125313}{Xie} & Long-term PM$_{2.5}$ and timing of pubertal development in boys and girls & \textit{Environ Res} & Prospective cohort & PM$_{2.5}$ & China\\
\addlinespace[0.9mm]
\multicolumn{6}{@{}l@{}}{\cellcolor{Mist}\textsf{\textbf{\textcolor{Sage}{Respiratory \& allergic}}}}\\[0.6mm]
\href{https://doi.org/10.1007/s00484-026-03289-z}{Bayar Muluk} & Climate change and otolaryngology: airway disease impacts and emission reduction & \textit{Int J Biometeorol} & Narrative review & PM$_{2.5}$ & Global\\
\href{https://doi.org/10.4103/mgr.MEDGASRES-D-25-00164}{Liao} & Short-term air pollution, respiratory symptoms and lung function in asthma & \textit{Med Gas Res} & Panel study & PM$_{2.5}$ / PM$_{10}$ & China\\
\href{https://doi.org/10.3390/epidemiologia7040097}{Sanchez de Prada} & Ambient PM and asthma-coded admissions in children 0-3 y, Spain & \textit{Epidemiologia} & Ecological panel & PM$_{2.5}$ / PM$_{10}$ & Spain\\
\href{https://doi.org/10.1016/j.jhazmat.2026.143076}{Yang} & Allergenic vs. toxic risk of airborne seed hairs compared with ambient PM$_{2.5}$ & \textit{J Hazard Mater} & Field measurement & PM$_{2.5}$ + bioaerosol & China\\
\addlinespace[0.9mm]
\multicolumn{6}{@{}l@{}}{\cellcolor{Mist}\textsf{\textbf{\textcolor{Violet}{Mechanistic toxicology}}}}\\[0.6mm]
\href{https://doi.org/10.1016/j.jes.2025.12.074}{Ding} & PM$_{2.5}$ promotes vascular calcification via miR-27b-5p/HES1 signalling & \textit{J Environ Sci (China)} & In vitro + organoid + murine & PM$_{2.5}$ & China\\
\href{https://doi.org/10.1007/s00360-026-01698-5}{Maraschi} & Osmoregulatory disruption and metal bioaccumulation from settleable atmospheric PM & \textit{J Comp Physiol B} & Controlled exposure (ecotox) & Settleable APM & Brazil\\
\href{https://doi.org/10.1016/j.jes.2025.11.018}{Wu} & Macrophage NLRP3 inflammasome in real-ambient PM-induced hepatic glucose dysregulation & \textit{J Environ Sci (China)} & Real-ambient murine + co-culture & Real-ambient PM & China\\
\addlinespace[0.9mm]
\multicolumn{6}{@{}l@{}}{\cellcolor{Mist}\textsf{\textbf{\textcolor{Clay}{Occupational \& indoor}}}}\\[0.6mm]
\href{https://doi.org/10.3389/fmedt.2026.1796828}{Chen} & Nested suction instrument to cut intraoperative surgical-smoke PM$_{2.5}$ & \textit{Front Med Technol} & Bench / device experiment & PM$_{2.5}$ (surgical smoke) & China\\
\href{https://doi.org/10.1093/annweh/wxag055}{Cole} & Determinants of respirable crystalline silica exposure in Australian tunnelling & \textit{Ann Work Expo Health} & Observational exposure study & Respirable dust / RCS & Australia\\
\href{https://doi.org/10.1016/j.jhazmat.2026.143081}{Ding} & Stage-resolved hair Zn-isotope dynamics as a firefighting combustion-exposure tracer & \textit{J Hazard Mater} & Repeated-measures biomarker & Combustion-derived PM & China\\
\href{https://doi.org/10.3390/jox16040119}{Ramsden} & Xenobiotic hazards in aircraft cabin air & \textit{J Xenobiot} & Review + risk economics & Ultrafine particles & Global\\
\addlinespace[0.9mm]
\multicolumn{6}{@{}l@{}}{\cellcolor{Mist}\textsf{\textbf{\textcolor{Amber}{Sensing, forecasting \& instrumentation}}}}\\[0.6mm]
\href{https://doi.org/10.1016/j.jenvman.2026.130515}{Shi} & Adaptive spatiotemporal graph transformer for multi-site multi-step PM$_{2.5}$ forecasting & \textit{J Environ Manage} & Deep-learning model & PM$_{2.5}$ & China / India\\
\href{https://doi.org/10.3390/bios16070380}{Tsuchiya} & Seasonal and indoor-outdoor characteristics of fluorescent aerosol particles in offices & \textit{Biosensors} & Real-time sensor campaign & Bioaerosol / coarse & Japan\\
\addlinespace[0.9mm]
\multicolumn{6}{@{}l@{}}{\cellcolor{Mist}\textsf{\textbf{\textcolor{Amber}{Exposure assessment \& modelling}}}}\\[0.6mm]
\href{https://doi.org/10.1029/2025GH001621}{Solorzano-Araque} & Multiscale exposure modelling: LOTOS-EUROS + mobility surveys + particle toxicity & \textit{GeoHealth} & CTM + mobility integration & PM$_{2.5}$ / PM$_{10}$ & Colombia\\
\addlinespace[0.9mm]
\multicolumn{6}{@{}l@{}}{\cellcolor{Mist}\textsf{\textbf{\textcolor{Coral}{Burden, policy \& mitigation}}}}\\[0.6mm]
\href{https://doi.org/10.1002/hsr2.72886}{Karamzad} & Ischaemic heart disease burden in MENA 1990-2021; PM attributable fraction 28.2\% & \textit{Health Sci Rep} & GBD secondary analysis & Ambient PM & MENA\\
\href{https://doi.org/10.1371/journal.pone.0347989}{Rees} & HEPA filtration to reduce respiratory infection in care homes & \textit{PLoS One} & RCT process evaluation & Filtered PM & UK\\
\href{https://doi.org/10.1016/j.jes.2026.01.009}{Shan} & Policy-driven emission reductions and post-pandemic rebound, high-resolution inventories & \textit{J Environ Sci (China)} & Emission inventory & PM emissions & China\\
\addlinespace[0.9mm]
\multicolumn{6}{@{}l@{}}{\cellcolor{Mist}\textsf{\textbf{\textcolor{Slate}{Other clinical endpoints}}}}\\[0.6mm]
\href{https://doi.org/10.3390/dj14070440}{Lupita} & PM$_{2.5}$, preventive-care capacity and caries/edentulism burden across the EU & \textit{Dent J} & Ecological panel & PM$_{2.5}$ & EU-27\\
\href{https://doi.org/10.1016/j.jdin.2026.06.008}{Park} & Long-term PM exposure and risk of androgenetic alopecia & \textit{JAAD Int} & Nationwide cohort & PM$_{2.5}$ / PM$_{10}$ & South Korea\\
\href{https://doi.org/10.3171/2026.3.PEDS2615}{Sun} & Environmental exposure and area deprivation vs. paediatric hydrocephalus outcomes & \textit{J Neurosurg Pediatr} & Registry cohort + geospatial & PM$_{2.5}$ (satellite-derived) & USA\\
\addlinespace[0.9mm]
\bottomrule
\end{longtable}}


\vspace{4mm}
\begin{keybox}
{\sffamily\bfseries\footnotesize\color{Deep}Provenance and method}\\[1.2mm]
{\fontsize{7.2}{8.8}\selectfont
Entry window \textbf{2026-08-17 to 2026-08-17}, contiguous with the 16 August issue; no gap.
\textbf{Per-source counts.} PubMed connector on two separate axes returned \textbf{9} on the
health axis and \textbf{1} on the sensing axis (the sensing hit was already in the health
set); the local \texttt{pubmed\_health} leg returned 7 and \texttt{pubmed\_sensing} 1, all
duplicates. Europe PMC on the primary particle query returned \textbf{1}, and a broader
air-quality sweep \textbf{6} of which 2 were new and in scope. Crossref-by-ISSN returned
\textbf{9 across 18 journals} and carried the whole instrumentation block --- \textit{Atmos
Chem Phys}, \textit{Aerosol Sci Technol}, \textit{J Aerosol Sci} and \textit{Aerosol Air
Qual Res} are not indexed in PubMed, so without this leg today's issue would have had no
sensing records at all. arXiv contributed \textbf{1} under a 3-day recency screen (a 14
August posting the 15 and 16 August sweeps both missed, which is the second time the
relevance query's unstable ordering has been caught by that screen).
\textbf{OpenAlex HTTP 429, 15th consecutive} --- the leg is dead and is no longer retried.
ClinicalTrials.gov returned \textbf{0} PM or air-pollution registrations created or updated
in the window; \texttt{trials.json} unchanged.\\[0.8mm]
\textbf{Two records are metadata-only.} Taylor \& Francis and Elsevier both withheld the
abstract for \texttt{10.1080/02786826.2026.2711663} and
\texttt{10.1016/j.jaerosci.2026.106883}, and the Semantic Scholar fallback was rate-limited
for the duration of the run. Both are labelled \texttt{Metadata only (no abstract)} in the
register, given their own slice in the architecture panel, and assert no finding beyond
their registered titles. Both are flagged for abstract retrieval on the next run. One record
is a \textbf{preprint} (Xia et al., arXiv) and is labelled as such.\\[0.8mm]
\textbf{Rejections (9).} Two out-of-scope PubMed hits (a black-carbon crop-physiology review
with no human endpoint or monitoring method; a family-medicine climate review naming air
pollution only as background); two Europe PMC drug-delivery records that matched on
\texttt{aerosol}/deposition terms but concern intranasal powder dosimetry and inhalable
antibacterial nanosystems; and five Crossref-by-ISSN records from in-scope journals that are
out of scope by subject --- aerosol--cloud effects on convective precipitation, urban heat
islands, Southern Ocean methylated sulfur, airborne eddy-covariance VOC fluxes, and an
opinion piece on AI in aerosol science. All logged with reasons to
\texttt{state/rejected.jsonl}.\\[0.8mm]
\textbf{Labels.} Three geography keys were added to \texttt{plots.py} this issue
(\texttt{Sweden}, \texttt{Finland}, \texttt{Pakistan}) plus four Swedish city needles;
\texttt{Sweden} had neither a key nor a needle and would have fallen to the unmapped bucket,
and \texttt{Pakistan} resolved to the correct group only by routing through an
\texttt{India} needle. No other new leaf label was written and the unmapped-label
diagnostics printed nothing. Assignment is single-label by dominant contribution and is a
judgement, not a controlled vocabulary. Every claim of the form "first/most/widest since X"
in the captions above was computed from \texttt{state/metrics.csv}, not recalled.
\texttt{check\_dois.py} returned \textbf{0 fail, 4 warn}; three warnings are title-overlap
heuristics firing on editorial focus lines and one is arXiv's absence from Crossref, each
confirmed by eye against the registered record.\\[0.8mm]
Content derived from PubMed, Europe PMC, Crossref and arXiv metadata; abstracts remain the
copyright of their respective publishers. All DOI links resolve to the publisher of record.\par}
\end{keybox}

\end{document}
